\documentclass[final]{article}


\usepackage{APS360}
\usepackage[utf8]{inputenc} % allow utf-8 input
\usepackage[T1]{fontenc}    % use 8-bit T1 fonts
\usepackage{hyperref}       % hyperlinks
\usepackage{url}            % simple URL typesetting
\usepackage{booktabs}       % professional-quality tables
\usepackage{amsfonts}       % blackboard math symbols
\usepackage{nicefrac}       % compact symbols for 1/2, etc.
\usepackage{microtype}      % microtypography
\usepackage{xcolor}         % colors
\title{Project Final Report for DriveVision}

\author{%
  Naz Saral \\
  1010184408, naz.saral@mail.utoronto.ca,
 \href{}{https://github.com/nazsaral5/project-aps360.git}
\\
}
\usepackage{graphicx}
\begin{document}
\maketitle
\vspace*{-0.5in}
\section*{Introduction}
\vspace*{-0.1in}
Traffic sign recognition systems play a fundamental role in developing technologies such as autonomous driving and driving assistance systems used in vehicles today. It is exciting to witness that vehicles today can interpret road rules and respond actively to their environment. This semester I worked on the traffic sign recognition tool called DriveVision. The goal of my project is to design and prototype an image classification tool that will recognize traffic signs around the globe, tested on our very own signs by the Ministry of Transportation (MTO) here in Ontario. I will scope down my project to identify five types of signage; the red and white octagonal stop sign, the upside down triangular yield sign, the rectangular and circular speed limit sign, the circular no-entry sign, and the rectangular one-way sign. This subset of signage influences driving, considering they directly indicate when to stop, speed up, slow down which is fundamental to the activity of operating a vehicle. The project will make use of deep learning concepts, as taught in this course, particularly for the training, validation, and testing of data. The dataset I will use to train the model is the The German Traffic Sign Recognition Benchmark (GTSRB), in addition to this I will also incorporate a collection of my own data from Ontario, including the rectangular speed limit and one-way signs to test the system. On the backend, the overall image classification will be done using convolutional neural networks (CNNs) to recognize and classify the signs. Model performance is evaluated using classification accuracy and confusion matrices on both the benchmark dataset and the Ontario evaluation dataset, a high-level illustration of DriveVision is captured in Figure 1.
\begin{figure}[!h]
  \centering
  \includegraphics[width=0.6\linewidth]{figure1.png}
  \caption{High-level overview of the DriveVision traffic sign recognition blackbox.}
  \label{fig:drivevision}
\end{figure}
\vspace{-0.3in}
\section*{Background & Related Work }
\vspace*{-0.1in}
For over 30 to 40 years researchers have iteratively developed feature classification systems for road signage based on edges, colors, shapes using image processing. Current research evolves within applications in daily driving such as autonomous vehicles, driver assistance systems and commercial vision pipelines. In recent years, related work that motivates this project for me are current technological advancements. In general traffic sign recognition is performed by feature matching, a first article mentions first preprocessing by converting RGB color space into HSV color space, then using transformation techniques to detect shapes (circles, triangles, squares), then applying feature matching methods to generate an output [1]. Another article explores the use of the GTSRB for traffic sign classification and recognition, using an improved LeNet-5 CNN model which is a highly complex network with seven layers, which include two convolutional layers, two pooling layers, two fully-connected layers and one output layer [2]. I will be drawing inspiration from both of these works to build out my model and establish a procedure as highlighted below for my DriveVision model. In addition to exploring the technical side of research in this field I also looked into specifics which motivate my own research such as Tesla Autopilot, where in the early 2020’s they updated their autopilot equipment to read speed limit signs and detect green lights being able to stop and slow down as needed, pioneering the commercialization of vision-based driver assistance systems for cars, initially launching the software back in 2019 [3]. Another popular application is Waymo’s self-driving technology, offering users in West Coast and Southern American states to enjoy a driverless rideshare. Benchmarking rideshare experiences with their own advancements in search and image recognition training, valuable in serving a public ridership [4]. Google Maps services make use of AI and imagery to identify speed limits, building on models trained on hundreds of different types of signs from all over the world [5]. All of these common applications of sign recognition explore the use of CNNs in detection algorithms as well as feature matching. 
\vspace*{-0.2in}
\section*{Data Processing}
\vspace*{-0.1in}
Upon collection of data from a German Traffic Sign Recognition Benchmark obtained from Kaggle. I began preprocessing data, as outlined in the proposal. Beginning with selecting the classes containing my 5 signs of focus; 0-8: 'speed\_limit', 13: 'yield', 14: 'stop', 17: 'no\_entry', 33-39: 'one\_way', and labelling them as such. I then counted and displayed the number of each type of sign (Figure 2), before assigning a 70/15/15 data split of training, validation and testing lending me to a split of; ​​Train: 15476, Val: 3316, Test: 3317. All images used were transformed with the following sequence of changes. 1. Scaling every image down to 32×32 using Resize(32, 32) effectively standardizes the input dimensions so the CNN always receives the same shape tensor. 2. ToTensor() to convert to tensor with shape [3, 32, 32] and rescale to range [0, 1]. 3. Then applied normalization, using ImageNet µ and σ values, which is standard for RGB channels. In addition to these changes, I implemented, 4. a random rotation of 15 decrees which will be effective especially with stop-signs and its octagonal shape, as well as, 5. colour adjustments to ensure random variation in the data which simulates different lighting conditions like overcast or sunny weather. Figure 3 is the output of processed subset. Some challenges presented in this phase show a noticeable class imbalance as the stop has only 780 samples vs. 12,780 for speed\_limit. Passing weights inversely proportional to class frequency into CrossEntropyLoss, so that mistakes on smaller classes are penalized more heavily will help with this challenge.
\begin{figure}[!h]
  \centering
  \includegraphics[width=0.8\linewidth]{figure23.png}
  \caption{Distribution of classes.}
  \caption{Sample of preprocessed images.}
  \label{fig:drivevision}
\end{figure}
\vspace*{-0.3in}
\section*{Architecture}
\vspace*{-0.1in}
\begin{figure}[!h]
  \centering
  \includegraphics[width=0.7\linewidth]{figure4.png}
  \caption{Pipeline Model of Architecture.}
  \label{fig:drivevision}
\end{figure}
For the image recognition, I have opted to use a Convolutional Neural Network setup, which will follow a sequential, layered, architecture that acts as an image processing pipeline. It will be capable of detecting edges, shapes as needed to classify the signs and decide which traffic sign it is.
A model as shown in Figure 4, that emulates what the  CNN will accomplish, is built with 3 convolutional layers, 3 pooling layers that are then flattened and passed through 2 fully connected layers. It contains nearly 620,000 trainable parameters. The validation accuracy of this model is nearly perfect at 0.999, as observed in Figure 5, the confusion matrix there is minimal error, with only 2/3317 samples of the most prevalent stop-sign class being misclassified. This is likely due to the quite low 32x32 resolution as seen in Figure 6.
\begin{figure}[!h]
  \centering
  \includegraphics[width=0.6\linewidth]{figure56.png}
  \caption{Confusion Matrix for Model.}
  \caption{Sample 32x32 input.}
  \label{fig:drivevision}
\end{figure}
\vspace*{-0.3in}
\section*{Baseline Model}
\begin{figure}[!h]
  \centering
  \includegraphics[width=0.6\linewidth]{figure7.png}
  \caption{Pipeline Model of Baseline Architecture.}
  \label{fig:drivevision}
\end{figure}
\vspace*{-0.2in}
\begin{figure}[!h]
  \centering
  \includegraphics[width=0.4\linewidth]{figure8.png}
  \caption{Confusion Matrix for Baseline Model.}
  \label{fig:drivevision}
\end{figure}
\vspace*{-0.1in}
The baseline model used is a Support Vector Machine (SVM). The structure of the model as shown in Figure 7, shows no spatial awareness, so just sees a flat list of pixel values. A quantitative measure of the success of this model is a 98.43\% validation accuracy, in addition to this, the confusion matrix (Figure 8) shows that the model performs the best on speed\_limit with 1917 samples being correctly identified. But a common error is that classes are occasionally misclassified as speed\_limit, this can be attributed to the fact it has the most samples. A challenge I faced training this model was a longer than usual runtime, which can be attributed to the large flattened inputs which is computationally expensive (roughly O(n²) to O(n³)) another issue at hand is that in future renditions, the model cannot generalize to Ontario signs as it is trained on European signs.
\vspace*{-0.2in}
\section*{Evaluation on New Data}
\vspace*{-0.1in}
\begin{figure}[!h]
  \centering
  \includegraphics[width=0.6\linewidth]{figure9.png}
  \caption{Processing of Training Images.}
  \label{fig:drivevision}
\end{figure}
\vspace*{0in}
In evaluating the model, previously the only images used were from the GTSRB, a collection of signs collected from all around Europe. As mentioned, I have been looking to implement local Ontario MTO signage to use on testing the DriveVision model, these images are a good representation of the model’s performance on new data. I went outside in my neighborhood and collected 25 images from around the Etobicoke area in order to create a test set. Of my signs of focus, I collected; 6 ‘speed\_limit’, 6 'yield', 3 'stop', 8 'no\_entry', 2 ‘one\_way.’ Images were processed following the process outlined above, as shown in Figure 9. The model was then evaluated on this data, to generate an accuracy of 0.480 as shown in the confusion matrix below, Figure 10.
\begin{figure}[!h]
  \centering
  \includegraphics[width=0.4\linewidth]{figure10.png}
  \caption{Confusion Matrix For Training.}
  \label{fig:drivevision}
\end{figure}
\vspace*{-0.3in}
\section*{Discussion}
\vspace*{-0.1in}
DriveVision performed best on the GTSRB, achieving an accuracy of 99.94\% versus the Baseline model, SVM's 98.43\%. The CNN clearly outperformed the baseline, which is expected considering its structure and complexity. The model was only troubled by two misclassifications out of 3317 samples and it converged quickly, reaching 99.97\% by the seventh epoch. When I tested the model on the curated Ontario dataset, it was interesting to see that accuracy dropped to 48\%, this is expected considering the model was trained on signs which differ in fonts, proportions, reflectivity. Once prevalent example of this is the speed limit sign, which is round in the dataset while being rectangular here in Ontario. Other reasons for this decrease in accuracy can be attributed to the resolution of images, noting that 32×32 may lose detail from photos taken at angles on my phone. Real-world conditions are also at hand here, where photos have trees, sky, shadows, damaged signage. Another interesting phenomenon is that a proportion of photos taken included stacked and layered signs where multiple signs were included in an image as can be seen in Figure 9, they were organized by their dominant sign, for example the one-way and no\_entry was placed into the no\_entry folder. It may also be useful to note that there was a class imbalance in the set of photos taken where there were only two one\_way and three stop samples compared to the six each in other categories. Overall, in using DriveVision, I learned that the 32×32 was too low for identifying features in an image and a higher resolution would help. Next, since the model trained only on European signs, as predicted it did not generalize well to Ontario’s signage, which could benefit from adding MTO signs to training data. 
\vspace*{-0.2in}
\section*{Ethical Considerations}
\vspace*{-0.1in}
As for ethical considerations for DriveVision, it is important to note that the tool’s goal is to decipher photos of signs into five possible categories. With such a task at hand some important considerations are as follows: \textbf{Misclassification:} This presents as a limitation to the model, considering traffic sign recognition is critical in applications of autonomous driving and driver’s assistance technology. If a sign is misclassified, there may pose risks to the driver and/or vehicle in operation while in traffic. It is important that in evaluating my model, that there is reasonable error, to ensure its safety.
\textbf{Dataset Representation and Generalization:} This is a limitation with respect to the data. Since the dataset I have chosen to use is primarily composed of European traffic signs. A model trained on one region may not generalize to another. Since I would like this model to be evaluated on Ontario traffic signs, which may differ in fonts, color schemes and layouts. To overcome this, I added a subset of relevant MTO signs that would be found within the province to strengthen my model's use here in Ontario. As this was only used in testing however, achieving an accuracy of 48\%, I know that DriveVision cannot be deployed safely for use in the province.
\textbf{Data Collection:} As a general ethical consideration, as mentioned, I will be including a subset of my own personal data collection in the sets. All traffic sign images collected by me were found within my neighbourhood in Etobicoke, Ontario, ensuring reliability and confidence.
\newpage
\section*{References}

[1] @inproceedings{ren2009traffic,
  author    = {Ren, Fei-Xiang and Huang, Jinsheng and Jiang, Ruyi and Klette, Reinhard},
  title     = {General Traffic Sign Recognition by Feature Matching},
  booktitle = {Proceedings of the 24th International Conference on Image and Vision Computing New Zealand},
  pages     = {409--414},
  year      = {2009},
  doi       = {10.1109/IVCNZ.2009.5378370}
}

[2] @article{cao2019traffic,
  title   = {Improved Traffic Sign Detection and Recognition Algorithm for Intelligent Vehicles},
  author  = {Cao, J. and Song, C. and Peng, S. and Xiao, F. and Song, S.},
  journal = {Sensors},
  volume  = {19},
  number  = {18},
  pages   = {4021},
  year    = {2019},
  doi     = {10.3390/s19184021}
}

[3] @online{bbc2020tesla,
  title        = {Tesla Updates Autopilot to Read Speed Signs},
  author       = {{BBC News}},
  organization = {BBC},
  year         = {2020},
  url          = {https://www.bbc.com/news/technology-53973511}
}

[4] @online{waymo2020perception,
  title        = {How Waymo Builds a Generalizable Perception System},
  author       = {{Waymo}},
  organization = {Waymo},
  year         = {2020},
  url          = {https://waymo.com/blog/2020/02/content-search}
}

[5] @online{googlemaps2020speedlimits,
  title        = {How AI and Imagery Keep Speed Limits on Google Maps Updated},
  author       = {{Google Maps}},
  organization = {Google},
  year         = {2020},
  url          = {https://blog.google/products-and-platforms/products/maps/how-ai-and-imagery-keep-speed-limits-on-google-maps-updated/}
}

[6] @misc{gtsrbkaggle,
  title        = {GTSRB German Traffic Sign},
  author       = {Meowmeowmeowmeowmeow},
  url = {https://www.kaggle.com/datasets/meowmeowmeowmeowmeow/gtsrb-german-traffic-sign}
}



\end{document}